% Options for packages loaded elsewhere
\PassOptionsToPackage{unicode}{hyperref}
\PassOptionsToPackage{hyphens}{url}
\PassOptionsToPackage{dvipsnames,svgnames,x11names}{xcolor}
\documentclass[
  american,
  11pt,
  letterpaper,
]{article}
\usepackage{xcolor}
\usepackage[margin=1in]{geometry}
\usepackage{amsmath,amssymb}
\setcounter{secnumdepth}{-\maxdimen} % remove section numbering
\usepackage{iftex}
\ifPDFTeX
  \usepackage[T1]{fontenc}
  \usepackage[utf8]{inputenc}
  \usepackage{textcomp} % provide euro and other symbols
\else % if luatex or xetex
  \usepackage{fontspec}
  \defaultfontfeatures{Scale=MatchLowercase}
  \defaultfontfeatures[\rmfamily]{Ligatures=TeX,Scale=1}
\fi
\usepackage{lmodern}
\ifPDFTeX\else
  % xetex/luatex font selection
  \setmainfont[]{Times New Roman}
\fi
% Use upquote if available, for straight quotes in verbatim environments
\IfFileExists{upquote.sty}{\usepackage{upquote}}{}
\IfFileExists{microtype.sty}{% use microtype if available
  \usepackage[]{microtype}
  \UseMicrotypeSet[protrusion]{basicmath} % disable protrusion for tt fonts
}{}
\makeatletter
\@ifundefined{KOMAClassName}{% if non-KOMA class
  \IfFileExists{parskip.sty}{%
    \usepackage{parskip}
  }{% else
    \setlength{\parindent}{0pt}
    \setlength{\parskip}{6pt plus 2pt minus 1pt}}
}{% if KOMA class
  \KOMAoptions{parskip=half}}
\makeatother
\usepackage{longtable,booktabs,array}
\newcounter{none} % for unnumbered tables
\usepackage{calc} % for calculating minipage widths
% Correct order of tables after \paragraph or \subparagraph
\usepackage{etoolbox}
\makeatletter
\patchcmd\longtable{\par}{\if@noskipsec\mbox{}\fi\par}{}{}
\makeatother
% Allow footnotes in longtable head/foot
\IfFileExists{footnotehyper.sty}{\usepackage{footnotehyper}}{\usepackage{footnote}}
\makesavenoteenv{longtable}
\ifLuaTeX
\usepackage[bidi=basic,shorthands=off]{babel}
\else
\usepackage[bidi=default,shorthands=off]{babel}
\fi
\ifPDFTeX
\else
\babelfont{rm}[]{Times New Roman}
\fi
\ifLuaTeX
  \usepackage{selnolig} % disable illegal ligatures
\fi
\setlength{\emergencystretch}{3em} % prevent overfull lines
\providecommand{\tightlist}{%
  \setlength{\itemsep}{0pt}\setlength{\parskip}{0pt}}
\usepackage{fancyhdr}
\usepackage{enumitem}
\usepackage{caption}
\usepackage{bm}
\usepackage{float}
\usepackage{microtype}
\usepackage{xurl}
\usepackage{ragged2e}

\pagestyle{fancy}
\fancyhf{}
\fancyfoot[C]{\thepage}
\renewcommand{\headrulewidth}{0pt}

\setlength{\parindent}{0pt}
\setlength{\parskip}{6pt plus 2pt minus 1pt}
\setlist{nosep,leftmargin=*}
\captionsetup{font=small,labelfont=bf}
\setlength{\emergencystretch}{3em}
\sloppy
\hyphenation{im-ple-men-ta-tion op-er-a-tion-al-i-za-tion}
\usepackage{bookmark}
\IfFileExists{xurl.sty}{\usepackage{xurl}}{} % add URL line breaks if available
\urlstyle{same}
\hypersetup{
  pdftitle={Cross-Disciplinary SIT Translation Guide},
  pdfauthor={Micah Blumberg},
  pdflang={en-US},
  colorlinks=true,
  linkcolor={blue},
  filecolor={Maroon},
  citecolor={blue},
  urlcolor={blue},
  pdfcreator={LaTeX via pandoc}}

\author{}
\date{}

\begin{document}

\begin{center}

{\LARGE\bfseries Cross-Disciplinary SIT Translation Guide\par}

\vspace{0.35em}

{\large\itshape Claim Labels, Evidence Roles, and Dated Status Boundaries\par}

\vspace{0.8em}

{\small\textbf{Author: Micah Blumberg}\par}

{\small Self Aware Networks Research Institute\par}

{\small\url{https://selfawarenetworks.com}\par}

{\small November 20, 2026 (Version 2 Portfolio Practical Index and Model Continuity Register to DOI 10.5281/zenodo.21693900)\par}

\end{center}

\vspace{0.65em}

\begin{abstract}

Research programs that connect physics, neuroscience, thermodynamics, computation, and mathematics face a recurring semantic hazard: a term can retain its spelling while changing its evidentiary status. \emph{Phase} may denote an angular field variable, a statistical relation between signals, or an informal stage. \emph{Projection} may denote a linear map, a measurement operation, a neural pathway, or a visualization. This guide supplies a controlled translation protocol for the Super Information Theory and Self-Aware Networks research program. Every cross-domain use receives a claim label, source class, domain limit, and demotion rule. Version 2 adds a dated six-axis status system that separates mathematical, empirical, visibility, package-validation, review, and release states. A generated appendix demonstrates the method on June and July 2026 proof, deposit, and revision events, including a public version 1 paired with a still-private version 2. The guide does not validate SIT, SAN, or an SIT / Riemann bridge. Its contribution is a reproducible governance method that prevents definitions from becoming data, analogies from becoming identities, computations from becoming theorems, and public deposits from becoming proof.

\end{abstract}

\section{Scope and evidence boundary}\label{scope-and-evidence-boundary}

As-of date for status examples: 2026-07-12\\
Public action authorized: no

\section{1. Reader contract}\label{reader-contract}

For any disputed sentence, ask four semantic questions and six status questions.

The semantic questions are:

\begin{enumerate}
\def\labelenumi{\arabic{enumi}.}
\tightlist
\item
  What kind of claim is being made?
\item
  What class of source supports that use?
\item
  What domain, model, population, theorem, or date bounds the source's authority?
\item
  How must the wording be demoted when it leaves that bound?
\end{enumerate}

The status questions are:

\begin{enumerate}
\def\labelenumi{\arabic{enumi}.}
\tightlist
\item
  What is proved or open mathematically?
\item
  What is tested empirically?
\item
  Is the artifact public or private?
\item
  Which declared package checks pass?
\item
  What kind of review has occurred?
\item
  Is release authorized?
\end{enumerate}

No answer may be inferred solely from another. This protocol is intentionally stricter than using one color or confidence score because a single scalar hides the reason a claim is limited.

A claim label and a status vector serve different purposes. The claim label identifies the kind of sentence being written. The status vector records the dated state of the exact claim and artifact. A sentence labeled ``formal theorem'' still needs scope, empirical, visibility, package, review, and release fields when those questions apply.

\section{2. Claim labels}\label{claim-labels}

\subsection{2.1 Definition and operational definition}\label{definition-and-operational-definition}

A definition fixes meaning inside a manuscript or model. An operational definition also specifies a measurement, calculation, schema, or decision rule. Neither is empirical confirmation. A variable named ``informational energy density'' remains a model quantity until a lawful measurement map and evidence connect it to conventional energy.

\subsection{2.2 Source-backed model term}\label{source-backed-model-term}

A local source can establish that an author defined and used a term. This supports exposition, chronology, and internal consistency. It does not establish conventional acceptance. SIT's local \texttt{R\_coh}, \texttt{J\_coh}, time-density, or coherence-current terms therefore require exact model equations and must not borrow the authority of standard currents or stress-energy tensors merely through similar notation.

\subsection{2.3 Proved-at-scope result}\label{proved-at-scope-result}

A proved result is identified by a statement, quantifiers, hypotheses, proof artifact, toolchain or review route, and exclusions. ``Lean checked theorem X'' is stronger than ``a chat proposed X'' but narrower than ``the intended research bridge is closed.'' Formal correctness and bridge adequacy are different review questions.

Four layers should accompany a load-bearing formal result: kernel and build status; the exact statement, quantifiers, and hypotheses; a mapping from source concepts to formal objects; and an independent review of whether the proved statement closes the intended bridge. A green first layer does not silently close the other three.

\subsection{2.4 Conditional theorem or yellow assumption-bound claim}\label{conditional-theorem-or-yellow-assumption-bound-claim}

A conditional result is valid only if its hypotheses hold. The manuscript must name the hypotheses near the claim. A proof of \texttt{H\ -\textgreater{}\ C} does not establish \texttt{H}. In cross-paper writing, an assumption may be imported only if the receiving paper preserves its open status.

\subsection{2.5 Diagnostic support}\label{diagnostic-support}

Scripts, simulations, benchmark tables, numerical fits, and traces can test implementation, illustrate behavior, discover counterexamples, or estimate parameters. They prove only claims about the computation unless a reviewed mathematical argument converts them into a theorem. A model fit also requires controls, uncertainty, and an empirical measurement map before it supports a physical claim.

\subsection{2.6 Empirical finding}\label{empirical-finding}

An empirical finding requires an identified experiment or observation, population and apparatus, analysis, uncertainty, and result. Compatibility with a mechanism is weaker than discrimination between mechanisms. Evidence for neural phase coding, traveling waves, or dendritic computation can support components of SAN without validating SAN's full conjunction or consciousness claim.

\subsection{2.7 Analogy, metaphor, and speculative bridge}\label{analogy-metaphor-and-speculative-bridge}

An analogy maps selected structure while preserving nonidentity. A metaphor organizes intuition and carries no inferential load. A speculative bridge proposes a relation whose formal or empirical link remains incomplete. These labels allow ambitious ideas to remain visible without presenting them as settled results.

\subsection{2.8 Prohibited promotion}\label{prohibited-promotion}

A prohibited promotion is not a forbidden topic. It is wording that exceeds the available support. Examples include converting a local package PASS into peer review, a finite theorem into an infinite operator result, a public archive into scientific correctness, or component neuroscience evidence into validation of a consciousness theory.

\section{3. Source classes and authority}\label{source-classes-and-authority}

\subsection{Primary literature}\label{primary-literature}

Peer-reviewed papers, scholarly books, standards, and primary datasets support conventional facts within their domains. Citation alone does not establish that an imported SIT or SAN interpretation follows from them.

\subsection{Formal artifacts}\label{formal-artifacts}

Proof-assistant files, checked kernels, theorem statements, claim signatures, and build certificates support exact formal claims. They require source identity, toolchain, build status, and the absence or presence of admitted statements. Formal artifacts do not supply empirical interpretation.

\subsection{Software and run outputs}\label{software-and-run-outputs}

Source code, tests, manifests, and deterministic outputs support reproducibility claims. Their authority is bounded by what the tests cover. A schema validator cannot certify physical truth; a numerical test cannot silently replace a proof obligation.

Package PASS is heterogeneous by design: it means that every check declared for that package passed. One package may rerun a model, another may validate a claim ledger, and another may verify a manifest. Comparisons therefore require each package's validation mode and key gate, not only the word PASS.

\subsection{Public archives and provider records}\label{public-archives-and-provider-records}

Zenodo, OSF, GitHub releases, and web archives establish availability, dates, files, checksums, and lineage when verified. They do not imply peer review. Provider state belongs in a visibility or release field, not a theorem field.

\subsection{Local sources and private provenance}\label{local-sources-and-private-provenance}

Drafts, notes, downloaded chats, and transcripts can recover chronology, terminology, donor documents, and research intent. They remain discovery or provenance evidence until a public-safe source is available. Absolute local paths belong in private manifests, not publication prose.

\subsection{AI-assisted material}\label{ai-assisted-material}

AI conversations and generated reviews can propose structures, objections, searches, or candidate citations. They are not primary scientific evidence and must not be used to certify novelty, equivalence, proof, priority, or empirical support. Their value is routing and adversarial drafting.

\section{4. Dated multidimensional status}\label{dated-multidimensional-status}

Program status is represented by six independent fields.

{\def\LTcaptype{none} % do not increment counter
\begingroup\footnotesize\setlength{\tabcolsep}{4pt}
\begin{longtable}[]{@{}
  >{\RaggedRight\arraybackslash}p{(\linewidth - 4\tabcolsep) * \real{0.3333}}
  >{\RaggedRight\arraybackslash}p{(\linewidth - 4\tabcolsep) * \real{0.3333}}
  >{\RaggedRight\arraybackslash}p{(\linewidth - 4\tabcolsep) * \real{0.3333}}@{}}
\toprule\noalign{}
\begin{minipage}[b]{\linewidth}\RaggedRight
Axis
\end{minipage} & \begin{minipage}[b]{\linewidth}\RaggedRight
Question
\end{minipage} & \begin{minipage}[b]{\linewidth}\RaggedRight
Example values
\end{minipage} \\
\midrule\noalign{}
\endhead
\bottomrule\noalign{}
\endlastfoot
mathematical & What statement is proved, conditional, open, or route-retired? & \texttt{proved-at-scope}, \texttt{conditional}, \texttt{open}, \texttt{retired-route} \\
empirical & What has been measured or tested? & \texttt{tested}, \texttt{diagnostic-only}, \texttt{not-tested} \\
visibility & Who can access the artifact? & \texttt{public}, \texttt{private}, \texttt{mixed} \\
package validation & Which declared integrity checks pass? & \texttt{pass}, \texttt{not-built}, \texttt{mixed} \\
review & What review or formal check occurred? & \texttt{formal-tool-check}, \texttt{internal-review}, \nolinkurl{independent-review-pending} \\
release & Is the exact artifact authorized and released? & \texttt{released}, \texttt{not-cleared}, \nolinkurl{blocked-pending-explicit-approval} \\
\end{longtable}
\endgroup
}

Every record also carries an as-of date, stable identity, exact sources, bounded wording, and a list of excluded promotions. Undated ``current'' prose is disallowed in generated status text because it cannot survive a later route or provider change.

In this table, \texttt{released} means that the identified version and file are publicly available through the named provider. It does not mean peer reviewed, independently reproduced, or scientifically endorsed. Mixed visibility must name the generations involved, such as public version 1 and private version 2, rather than implying unexplained provider disagreement.

\section{5. Physics vocabulary}\label{physics-vocabulary}

\subsection{Coherence}\label{coherence}

In standard physics, coherence may describe stable phase relations, off-diagonal density-matrix structure, interference visibility, or an order parameter under specified conditions. In SIT, coherence can also be a model-organizing term. Safe writing identifies the mathematical object and its units or normalization. ``Coherence'' alone does not establish superconductivity, quantum coherence, global synchronization, low entropy, or consciousness.

\subsection{Phase and phase locking}\label{phase-and-phase-locking}

Phase may be a gauge-dependent field coordinate, a gauge-invariant difference or holonomy, or an oscillatory signal relation. A physical claim should identify the observable invariant. Phase locking requires a dynamics, coupling, state space, and lock criterion. STP / GPL language remains a model proposal unless a gravity-to-parameter map and empirical protocol are supplied.

\subsection{Projection}\label{projection}

Projection can denote an idempotent operator, a map from one state space to another, a measurement operation, or an informal rendering. A theorem about an abstract finite edge-fiber kernel does not establish route-specific physical readout. The receiving manuscript must repeat the hypotheses and state whether the map is mathematical, computational, or measured.

\subsection{Master action and effective-field-theory language}\label{master-action-and-effective-field-theory-language}

An action can organize fields and couplings inside a proposed model. Calling it EFT-like requires a declared field content, symmetries, operator basis, dimensions, cutoff, and approximation regime. It does not imply ultraviolet completion or empirical adequacy. Noether-current language requires the symmetry assumptions and the open-system caveat when dissipation is present.

\section{6. Neuroscience and SAN vocabulary}\label{neuroscience-and-san-vocabulary}

\subsection{Tonic and phasic}\label{tonic-and-phasic}

Tonic and phasic should be treated as role-relative dynamical descriptions unless a study fixes them to particular cells or frequency bands. A tonic process supplies a comparatively persistent baseline over a declared window; a phasic process supplies a transient deviation relative to that baseline. This avoids turning delta, theta, alpha, beta, or gamma into permanent functional labels.

\subsection{Phase coding and phase-wave differential}\label{phase-coding-and-phase-wave-differential}

Standard phase coding concerns relationships between neural events and an oscillatory reference. The SAN phase-wave differential is a framework variable and requires a receiver, reference, observable, and scale. Evidence for phase coding or traveling waves supports narrower ingredients; it does not by itself establish the SAN variable or a global electromagnetic phase landscape.

\subsection{Neural rendering and NAPOT}\label{neural-rendering-and-napot}

Neural rendering and Neural Array Projection Oscillation Tomography describe a proposed recurrent, multiscale reconstruction process. ``Tomography'' is an analogy or model term unless an inverse problem, forward operator, identifiability conditions, and reconstruction validation are specified. The phrase ``3D volumetric television'' is explanatory metaphor, not evidence.

\subsection{Dendritic token and broadcast}\label{dendritic-token-and-broadcast}

A dendritic branch can perform nonlinear, state-dependent integration. Calling its output a token is safe only when the manuscript defines the compressed state variables and downstream decoding. Synaptic, axonal, electrical-field, and ephaptic effects must remain distinct propagation terms. Ephaptic coupling can be discussed as a local contribution where supported; it should not be promoted to a sole global broadcast mechanism.

\subsection{Consciousness and self-model claims}\label{consciousness-and-self-model-claims}

Evidence for prediction, reafference, body maps, traveling waves, synchrony, or dendritic computation does not establish consciousness sufficiency. SAN's conjunction should be stated as a testable synthesis with discriminating experiments and strong alternatives.

\section{7. Thermodynamics and dissipation vocabulary}\label{thermodynamics-and-dissipation-vocabulary}

Thermodynamic entropy, information-theoretic entropy, free energy, prediction error, signal decay, and informal disorder are not interchangeable. Every equation must state the quantity, units, ensemble or distribution, and system boundary.

``Dissipation'' can mean irreversible conversion of free energy, damping of a signal, relaxation of a model variable, or absorption of mismatch into an updated predictive state. The last is a proposed computational interpretation. It should not be called a new thermodynamic law without a state function, dynamics, relation to established laws, limiting cases, and empirical tests.

\section{8. Computation and proof vocabulary}\label{computation-and-proof-vocabulary}

A computational witness is an object checked by a verifier. That vocabulary does not prove a complexity-class containment, separation, or collapse. P versus NP claims require standard language definitions, quantified complexity bounds, reductions, and proof. Search heuristics and observed runtimes remain diagnostic.

Similarly, ``proof packet'' can mean an organized evidence bundle rather than a mathematical proof. The guide reserves \texttt{proved-at-scope} for an exact result with a valid proof route. Manifests, checksums, and tests establish artifact integrity and reproducibility; they do not upgrade an open theorem.

\section{9. Translation procedure}\label{translation-procedure}

For each cross-domain sentence:

\begin{enumerate}
\def\labelenumi{\arabic{enumi}.}
\tightlist
\item
  Identify the native-domain meaning of the key term.
\item
  State the receiving-domain meaning.
\item
  Name the invariant structure preserved by the transfer.
\item
  Name what is not preserved.
\item
  Assign claim and source labels.
\item
  Add the applicable status vector and as-of date.
\item
  State a falsifier, discriminator, or reopen condition when the claim is substantive.
\end{enumerate}

For example, a neural oscillation and a field-theory phase both admit angular descriptions, but the objects, symmetries, observables, scales, and evidence differ. The shared mathematics can motivate a model. It does not establish physical identity.

\section{10. Demotion rules}\label{demotion-rules}

{\def\LTcaptype{none} % do not increment counter
\begingroup\footnotesize\setlength{\tabcolsep}{4pt}
\begin{longtable}[]{@{}
  >{\RaggedRight\arraybackslash}p{(\linewidth - 2\tabcolsep) * \real{0.5000}}
  >{\RaggedRight\arraybackslash}p{(\linewidth - 2\tabcolsep) * \real{0.5000}}@{}}
\toprule\noalign{}
\begin{minipage}[b]{\linewidth}\RaggedRight
Risk
\end{minipage} & \begin{minipage}[b]{\linewidth}\RaggedRight
Required correction
\end{minipage} \\
\midrule\noalign{}
\endhead
\bottomrule\noalign{}
\endlastfoot
analogy written as identity & Name the shared structure and the nonidentity. \\
model definition written as measurement & Add the measurement map or label the quantity model-defined. \\
finite result promoted to infinite theorem & Restore quantifiers and list the missing limit or operator obligations. \\
diagnostic promoted to proof & Relabel as computation and identify the unresolved proof obligation. \\
component evidence promoted to framework validation & Name the supported component and the untested conjunction. \\
public deposit promoted to peer review & Separate visibility, review, and claim status. \\
package PASS promoted to scientific acceptance & List the checks covered and retain scientific / review gates. \\
retired route generalized to all routes & Name the route family and decision date. \\
private path used as public citation & Replace with a stable public source or retain as private provenance only. \\
AI output used as authority & Use it for discovery; verify against primary or formal sources. \\
\end{longtable}
\endgroup
}

\section{11. Dated case study}\label{dated-case-study}

The generated appendix records eight examples. The June 1 ledgers accept a bounded model surface and conditional finite kernel while disposing named route families and leaving full SIT and RH open. BC-055O later proves one finite constructor contract without reopening those higher claims. The July 7 deposit wave changes ten records from private to public without changing theorem status. The pre-rank-28-registration July 12 aggregate shows 32 passing private packages and zero public- action authorizations. The STP companion combines proved model results, no empirical GPL test, a passing package, pending independent review, and blocked release. The guide's own mixed-visibility row names a public version 1 and a private version 2; it is not a provider-parity claim.

These records would be contradictory in a one-dimensional ladder. They are coherent in a status vector because they answer different questions.

\section{12. Reproducibility and provenance}\label{reproducibility-and-provenance}

Ten source files are frozen in the private revision package. Each manifest row records the package- local snapshot, original disk path, byte count, SHA-256 digest, and role. The status generator fails on source drift, missing dates, unknown labels, missing BC-055O exclusions, private-path leakage, visibility / proof label substitution, or generated-output drift.

Automated path rejection is not a complete privacy review. Before release, every file must be classified as private-only, public-candidate, or generated-public-candidate and checked for names, unreleased identifiers, source excerpts, rights, and absolute paths. This package remains private.

The public manuscript should cite stable provider URLs, DOIs, standards, and formal artifact releases. The private evidence package may retain exact disk paths to support audit and restart. This separation protects privacy without discarding provenance.

\section{13. Conventional anchors}\label{conventional-anchors}

The guide uses conventional sources to define standard vocabulary, not to validate SIT or SAN:

\begin{itemize}
\tightlist
\item
  Aharonov and Bohm, ``Significance of Electromagnetic Potentials in the Quantum Theory,'' \emph{Physical Review} 115, 485-491 (1959), DOI \nolinkurl{10.1103/PhysRev.115.485}.
\item
  Kuramoto, \emph{Chemical Oscillations, Waves, and Turbulence} (1984), DOI \nolinkurl{10.1007/978-3-642-69689-3}.
\item
  Fries, ``Rhythms for Cognition: Communication through Coherence,'' \emph{Neuron} 88, 220-235 (2015), DOI \nolinkurl{10.1016/j.neuron.2015.09.034}.
\item
  O'Keefe and Recce, ``Phase relationship between hippocampal place units and the EEG theta rhythm,'' \emph{Hippocampus} 3, 317-330 (1993), DOI \nolinkurl{10.1002/hipo.450030307}.
\item
  Rao and Ballard, ``Predictive coding in the visual cortex,'' \emph{Nature Neuroscience} 2, 79-87 (1999), DOI \texttt{10.1038/4580}.
\item
  W3C, ``PROV-O: The PROV Ontology'' (2013), \url{https://www.w3.org/TR/prov-o/}.
\item
  FORCE11, ``Joint Declaration of Data Citation Principles'' (2014), DOI \nolinkurl{10.25490/a97f-egyk}.
\end{itemize}

\section{14. Limitations}\label{limitations}

This guide cannot adjudicate the correctness of every imported domain claim. Its source set is a dated program snapshot, not a complete literature review. Its generated checks detect declared status and provenance errors, not all scientific errors. Independent domain review remains necessary, especially for mathematical scope, field-theory language, neuroscience causality, and thermodynamic interpretation.

The method also imposes documentation cost. That cost is justified for load-bearing claims but can be reduced for ordinary prose. The appropriate unit is the material inference, not every noun.

The status appendix is a historical snapshot, not a live dashboard. A later program change should produce a new dated snapshot and source manifest rather than silently rewriting this one.

\section{15. Review and release boundary}\label{review-and-release-boundary}

The schema, generator, and internal adversarial pass are complete for this near-final draft. Independent domain review remains open for mathematical scope, source interpretation, physics, neuroscience, and claim language. Micah Blumberg must approve the exact final Markdown, PDF, metadata, and repository releases. No Zenodo, OSF, GitHub, DOI, journal, or other public mutation is authorized by this draft.

\section{Conclusion}\label{conclusion}

Cross-disciplinary rigor requires two kinds of translation: semantic translation between domains and temporal translation between program states. The first asks what a term means and what evidence supports it. The second asks what changed, when, and on which independent status axis. Version 2 makes both explicit. Its governing rule is simple: define the claim, identify the source, bound the domain, date the state, preserve excluded promotions, and never let one status field impersonate another.

\section{Appendix A: Dated Status Snapshot}\label{appendix-a-dated-status-snapshot}

As-of date: 2026-07-12. This appendix separates mathematical, empirical, visibility, package-validation, review, and release status. No column determines another.

Frozen source set: 10 files. Pre-rank-28-registration package snapshot: 32/32 PASS, with 0 public- action authorizations. Branch C provider snapshot: 10 Zenodo-open records mirrored on public OSF node \texttt{w5uy7}.

{\def\LTcaptype{none} % do not increment counter
\begingroup\tiny\setlength{\tabcolsep}{2pt}
\begin{longtable}[]{@{}
  >{\RaggedRight\arraybackslash}p{(\linewidth - 16\tabcolsep) * \real{0.1111}}
  >{\RaggedRight\arraybackslash}p{(\linewidth - 16\tabcolsep) * \real{0.1111}}
  >{\RaggedRight\arraybackslash}p{(\linewidth - 16\tabcolsep) * \real{0.1111}}
  >{\RaggedRight\arraybackslash}p{(\linewidth - 16\tabcolsep) * \real{0.1111}}
  >{\RaggedRight\arraybackslash}p{(\linewidth - 16\tabcolsep) * \real{0.1111}}
  >{\RaggedRight\arraybackslash}p{(\linewidth - 16\tabcolsep) * \real{0.1111}}
  >{\RaggedRight\arraybackslash}p{(\linewidth - 16\tabcolsep) * \real{0.1111}}
  >{\RaggedRight\arraybackslash}p{(\linewidth - 16\tabcolsep) * \real{0.1111}}
  >{\RaggedRight\arraybackslash}p{(\linewidth - 16\tabcolsep) * \real{0.1111}}@{}}
\toprule\noalign{}
\begin{minipage}[b]{\linewidth}\RaggedRight
Record
\end{minipage} & \begin{minipage}[b]{\linewidth}\RaggedRight
As of
\end{minipage} & \begin{minipage}[b]{\linewidth}\RaggedRight
Mathematical
\end{minipage} & \begin{minipage}[b]{\linewidth}\RaggedRight
Empirical
\end{minipage} & \begin{minipage}[b]{\linewidth}\RaggedRight
Visibility
\end{minipage} & \begin{minipage}[b]{\linewidth}\RaggedRight
Package
\end{minipage} & \begin{minipage}[b]{\linewidth}\RaggedRight
Review
\end{minipage} & \begin{minipage}[b]{\linewidth}\RaggedRight
Release
\end{minipage} & \begin{minipage}[b]{\linewidth}\RaggedRight
Bounded wording
\end{minipage} \\
\midrule\noalign{}
\endhead
\bottomrule\noalign{}
\endlastfoot
\nolinkurl{SIT-CAPPED-BASELINE-20260601} & 2026-06-01 & mixed-bounded & not-applicable & private & mixed & internal-review & not-cleared & The 2026-06-01 ledgers accept the named model surface and conditional finite kernel at their stated scopes. \\
\nolinkurl{SIT-ROUTES-DISPOSED-20260601} & 2026-06-01 & retired-route & not-applicable & private & mixed & internal-review & not-cleared & The three named route families were retired as unconditional routes in the 2026-06-01 proof-stack review. \\
\nolinkurl{SIT-RIEMANN-CAP-20260601} & 2026-06-01 & open & not-applicable & private & mixed & internal-review & not-cleared & The 2026-06-01 ledger retires its named finite symbolic route without proving or disproving the Riemann Hypothesis. \\
\nolinkurl{BC-055O-CONSTRUCTOR-CLOSURE-20260616} & 2026-06-16 & proved-at-scope & not-applicable & private & pass & formal-tool-check & not-cleared & BC-055O closes its named finite Farey constructor contract in Lean. \\
\nolinkurl{BRANCH-C-PUBLIC-DEPOSITS-20260707} & 2026-07-07 & mixed-bounded & mixed & public & pass & mixed & released & The ten-paper Branch C series is public, while each mathematical claim remains limited by its own proof and review status. \\
\nolinkurl{REVISION-AGGREGATE-20260712} & 2026-07-12 & not-applicable & not-applicable & private & pass & internal-review & blocked-pending-explicit-approval & The pre-rank-28-registration package snapshot passes its declared structural and computational checks; scientific acceptance and release approval remain separate. \\
\nolinkurl{STP-FINITE-LOCK-COMPANION-20260712} & 2026-07-12 & proved-at-scope & not-tested & private & pass & independent-review-pending & blocked-pending-explicit-approval & The private companion proves bounded model results; its physical interpretation, independent review, and release remain open. \\
\nolinkurl{TRANSLATION-GUIDE-V2-20260712} & 2026-07-12 & not-applicable & not-applicable & mixed & not-built & independent-review-pending & blocked-pending-explicit-approval & Version 1 is public, while version 2 remains a private revision until validation, review, and explicit author approval are complete. \\
\end{longtable}
\endgroup
}

\subsection{Excluded promotions}\label{excluded-promotions}

\begin{itemize}
\tightlist
\item
  \nolinkurl{SIT-CAPPED-BASELINE-20260601}: route-specific projection or readout; Coherence Tower construction; physical SIT; full SIT; SIT / Riemann; Riemann Hypothesis; publication clearance.
\item
  \nolinkurl{SIT-ROUTES-DISPOSED-20260601}: all later finite routes are retired; no later subroute can close; full SIT is disproved; Riemann Hypothesis is resolved.
\item
  \nolinkurl{SIT-RIEMANN-CAP-20260601}: SIT proves the Riemann Hypothesis; the Riemann Hypothesis is disproved; all Branch C finite results are invalid.
\item
  \nolinkurl{BC-055O-CONSTRUCTOR-CLOSURE-20260616}: finite C-22; infinite transfer; SIT / Riemann bridge; HSIT Riemann; Riemann Hypothesis; full SIT; physical interpretation.
\item
  \nolinkurl{BRANCH-C-PUBLIC-DEPOSITS-20260707}: public means peer reviewed; public means independently reproduced; the series proves the Riemann Hypothesis; the series proves full SIT.
\item
  \nolinkurl{REVISION-AGGREGATE-20260712}: package PASS proves scientific correctness; package PASS is independent review; package PASS authorizes publication.
\item
  \nolinkurl{STP-FINITE-LOCK-COMPANION-20260712}: gravity-to-parameter derivation; empirical GPL validation; Bell refutation; full SIT proof; release authorization.
\item
  \nolinkurl{TRANSLATION-GUIDE-V2-20260712}: revision already released; package already approved; provider mutation authorized.
\end{itemize}

\subsection{Interpretation rule}\label{interpretation-rule}

A public artifact can contain an open claim. A checked theorem can remain private. A package can pass its declared checks while independent review and release authorization remain pending. These combinations are expected states, not contradictions.

Historical July 2026 release status preserved above; November 20, 2026 private review candidate status recorded in Appendix B.4.

\section{Appendix B: November 20, 2026 Portfolio Practical Index, Model Continuity, and Version Lineage}\label{appendix-b-november-20-2026-portfolio-practical-index-model-continuity-and-version-lineage}

This appendix updates the translation guide into a practical claim-to-model index across the eleven upgraded intervening SIT papers (following the November 20, 2026 scientific correction pass), the version-aware neuroscience track (Paper 114), and the separate SIT Version 81 monograph lane. Every material claim carries its owning model family, core assumptions, exact verification script, unresolved dependency, and public DOI lineage.

\subsection{B.1 Corrected Version Chronology}\label{b.1-corrected-version-chronology}

\begin{enumerate}
\def\labelenumi{\arabic{enumi}.}
\tightlist
\item
  \textbf{SIT Version 79 (December 18, 2025):} Omnibus monograph baseline (DOI 10.5281/zenodo.18133332).
\item
  \textbf{July 2026 Version 80 Compact Corrective Lineage:} Version 80 Draft 13 (July 26, 2026; frozen comparison source for Paper 048) released publicly on July 29, 2026 as \textbf{Paper 002} (DOI 10.5281/zenodo.21693867).
\item
  \textbf{November 2026 Formal Monograph Series (renamed SIT Version 81):} Stage 5.2 (v80.2, November 16, 2026) and Stage 6 (v80.3, commits \nolinkurl{30a71e4ffdda366c8d2489ee553d4a0a8b4c17af} and \texttt{d5dc6412da}, November 18--19, 2026), renamed \textbf{SIT Version 81} on its independent monograph branch (\nolinkurl{codex/sitproof-sit81-monograph-20261119}) to preserve chronological clarity relative to July Version 80.
\item
  \textbf{Paper 114 (\emph{Learning to Feel}) Version-Aware Record:} Initial public release \texttt{v1.0} on October 17, 2026 (DOI 10.5281/zenodo.23447925); current main paper \texttt{v1.2} dated October 20, 2026 (\url{https://pub.svgn.io/papers/paper-114/}); technical supplement deliberately retains its \texttt{v1.1} label as documented version history rather than a naming defect. Paper 114 continues on its independent neuroscience track with \textbf{no identification} between neural synchrony/learned influence and canonical SIT scalar fields.
\end{enumerate}

\subsection{B.2 Unified Clock and Observable Dictionary (Explicit GR Baseline Separation)}\label{b.2-unified-clock-and-observable-dictionary-explicit-gr-baseline-separation}

Across all portfolio papers: \(r := \delta\rho_t/\rho_0\) is the dimensionless field perturbation; \(\alpha_\rho\) is the local intrinsic (fixed-metric) clock-response coefficient (\(y_{\mathrm{clock,int}} = \alpha_\rho r\)); \(\mathcal{T}\) is the gravitational transfer (\(r = \mathcal{T} p\), \(p = \Delta\Phi/c^2\)); and the total observed weak-field gravitational redshift slope is \(d\ln\nu / d(\Phi/c^2) = 1 + \delta\alpha_\Phi\), where \(1\) is the standard General Relativity metric redshift and \(\delta\alpha_\Phi := \alpha_\rho\mathcal{T}\) is the \textbf{anomalous (non-metric, SIT-additional) fractional redshift slope} in which \(\mathcal{T}\) enters \textbf{strictly once}. Point-process event rates \(N(\tau)/\tau\) (Paper 007) and atomic clock frequency shifts \(\Delta\nu/\nu\) are distinct operational observables requiring calibrated detector models.

\subsection{B.3 Master Claim-Model-Evidence Index (November 20, 2026 Corrected Editions)}\label{b.3-master-claim-model-evidence-index-november-20-2026-corrected-editions}

{\def\LTcaptype{none} % do not increment counter
\begingroup\footnotesize\setlength{\tabcolsep}{4pt}
\begin{longtable}[]{@{}
  >{\RaggedRight\arraybackslash}p{(\linewidth - 6\tabcolsep) * \real{0.2500}}
  >{\RaggedRight\arraybackslash}p{(\linewidth - 6\tabcolsep) * \real{0.2500}}
  >{\RaggedRight\arraybackslash}p{(\linewidth - 6\tabcolsep) * \real{0.2500}}
  >{\RaggedRight\arraybackslash}p{(\linewidth - 6\tabcolsep) * \real{0.2500}}@{}}
\toprule\noalign{}
\begin{minipage}[b]{\linewidth}\RaggedRight
Paper \& Public DOI
\end{minipage} & \begin{minipage}[b]{\linewidth}\RaggedRight
Model Family
\end{minipage} & \begin{minipage}[b]{\linewidth}\RaggedRight
Verified November 2026 Upgrade \& Script
\end{minipage} & \begin{minipage}[b]{\linewidth}\RaggedRight
Unresolved Dependency / Scope Ceiling
\end{minipage} \\
\midrule\noalign{}
\endhead
\bottomrule\noalign{}
\endlastfoot
\textbf{Paper 002} (\nolinkurl{10.5281/zenodo.21693867}) & \nolinkurl{M_JUL2026_CONFORMAL} + crosswalk to \nolinkurl{M_NOV2026_AUTH_EFT_V1} & States exact non-Weyl-invariant matter assumptions for conformal obstruction (\(A = f_{\mathrm{SM}} = 1\)), notes total \(\mathcal{L}_m = 0\) for free/bilinear Dirac, qualifies visible-gauge photon extension, adopts \(1 + \delta\alpha_\Phi\) clock dictionary, and locks Reference Cases A/B/C (\nolinkurl{verify_paper002_v2.py}, 4/4 pass) & Microscopic composite-body charges, visible \(U(1)\) photon dictionary, and quantum continuum limit remain open \\
\textbf{Paper 003} (\nolinkurl{10.5281/zenodo.21693868}) & \nolinkurl{M_LEVEL0_CONDENSATE} (distinct phenomenological model) & Reconciles canonical scalar source \((\nabla^2 - m_\sigma^2)\delta\sigma_Q = \lambda_Q n_0\) (\(\lambda_Q := g_Q m_a / M_{\mathrm{Pl}}\)) and \(g_Q = 0\) null condition; separates conservative GP (\(\Gamma_{12}=0\)) from open ZNG-type transfer (\(-i\hbar\Gamma_{12}\Phi_0/(2n_0)\) + paired thermal continuity), exact full-potential energy balance \(n_0\partial_t(V_{\mathrm{trap}}+U)\), and \(^{87}\mathrm{Rb}\) confounder budget (\nolinkurl{verify_paper003_v4.py}, 3/3 pass) & Nonlocal Penrose--Onsager localization valid only in coarse hydrodynamic frame; dynamical scalar-field energy exchange, full ZNG kinetic collision model, \& covariant completion open \\
\textbf{Paper 007} (\nolinkurl{10.5281/zenodo.21693874}) & \nolinkurl{M_SDT_ONE_SCALAR} (static Schur vs.~dynamical reduction) & Distinguishes static Schur \(m_{\mathrm{eff}}^2\) from dynamical eigenvalue (\(0.75\) vs.~\(0.6972\)), states \(\lambda_{\mathrm{mix}}=J_R=0\) freezing condition, adds exact renewal Fano \(F_{\mathrm{ren}} = 1/(1+\eta_{\mathrm{occ}})^2 = 0.25\) alongside assumed \(F=1, 2.5\) (\nolinkurl{verify_paper007_v2.py}, 2/2 pass) & Microscopic quantum-event-to-gravity derivation remains open; scalar elimination does not equate conformal and multiplicative matter actions \\
\textbf{Paper 014} (\nolinkurl{10.5281/zenodo.21693886}) & \nolinkurl{M_LEVEL0_CONDENSATE} synthetic \(^{87}\mathrm{Rb}\) interferometer & Synthetic conditional FIM/VIF design (\(7.14\times\) sharper); exact multivariable EIV attenuation (\(38.09\%\) bias in collinear Design ii vs.~\(1.19\%\) in orthogonal Design iii), common-additive-offset slope invariance vs.~classical/correlated EIV hierarchy, and scoped negative reach for declared \(10^7\)-atom BEC source (\nolinkurl{verify_paper014_v3.py}, 2/2 pass) & Actual laboratory calibration records, hierarchical calibration/error model, and Section 6 material-charge/screening translation for any condensate-selective torsion bound \\
\textbf{Paper 028} (\nolinkurl{10.5281/zenodo.21693900}) & Portfolio governance \& translation index & Unites chronology, GR-separated clock dictionary, all 11 corrected SIT upgrades, Paper 114 version-aware entry, and SIT 81 cross-reference & Living index; updated whenever an owning paper advances \\
\textbf{Paper 033} (\nolinkurl{10.5281/zenodo.21693908}) & \nolinkurl{M_LATTICE_U1_2D_TOY} & Fixes Thm 6 \(f_j \mid f_{j+1}\) divisibility, RD3.1 \& Thm 10 Step 3 fibre multiplicity; proves Thm 12 (\(256 \to 136\) Weyl gauge quotient vs.~exact \(136\)-class bijection from larger Wilson loops) \& Thm 13 (exact necessary (19a) vs.~sufficient (19b) under full word realizability, fixed-\(K_1\) polynomial feasibility \(1+\gamma < 2K_1\) with \(m=29\) vs.~\(\gamma=3\) impossibility, and \(\delta_{\mathrm{margin}}>0\) noisy Remark R13.2) (\nolinkurl{verify_paper033_v2.py}, 4/4 pass) & Non-abelian \(SU(2)\) adapter (either accepting 136-class quotient or proving matter-enlarged relational separation), 4D Bianchi constraints, and continuum Haar measure remain open \\
\textbf{Paper 034} (\nolinkurl{10.5281/zenodo.21693909}) & Conditional measurement-dependent phase-lock model & Establishes Wilson-line gauge-invariant relative phase requirement; proves \& verifies Thm 6A.1 transient Fokker--Planck \(\pi\)-symmetry (\nolinkurl{verify_paper034_v2.py}, 2/2 pass) & Causal gravitational preparation mechanism, losses, and timing loopholes remain open \\
\textbf{Paper 048} (\nolinkurl{10.5281/zenodo.21693926}) & July v80 Draft 13 + November \nolinkurl{M_NOV2026_AUTH_EFT_V1} & Verifies exact canonical November linked-square potential \& mixed Hessian \(H_{12} \neq 0\); derives exact interacting Dirac on-shell law \(\mathcal{L}_m = (1-n)\mathcal{L}_{\mathrm{int}}\) (\(n=1\): \(\mathcal{L}_{\mathrm{Dirac}}=-\mathcal{L}_{\mathrm{int}}, \mathcal{L}_m=0\); \(n=2\): \(-f'_0\mathcal{L}_{\mathrm{int}}\)) and explicit dressed \(f_0 q_\Psi \bar{\Psi}\gamma^\mu\Psi\) vs.~canonical \(\tilde{\Psi}\) gauge-current normalization (\nolinkurl{verify_paper048_v2.py}, 1/1 pass) & Shared constant-background free/bilinear limit retains strictly zero discriminatory value between CFS and SIT \\
\textbf{Paper 055} (\nolinkurl{10.5281/zenodo.21693935}) & Standing counterexample gate library & Verifies 4 explicit countermodels: non-unique kernels (\(\mu_4\) ratio \(3\), \(\kappa_4 = -2h^4\) vs.~\(0\)), overdamped drift-velocity contrast \(\Delta_{v,AB}=2/3\) with \(d^2\langle x\rangle/dt^2=0\), undamped heating runaway, and clock-vs-deflection underdetermination (\nolinkurl{verify_paper055_v2.py}, 4/4 pass) & Any future QGTCD/SIT derivation must state derived vs.~assumed premise excluding each countermodel \\
\textbf{Paper 062} (\nolinkurl{10.5281/zenodo.21693943}) & Classical QGTCD vs.~quantum entropic bath & Derives exact conditional thermal calibration \(c_A = 2\beta k_B T\), corrected amplitude decoherence \(\Gamma_{\mathrm{dec}} = k_B T (\Delta x)^2/(\mu\hbar^2)\) (\(3.43\times 10^8\) to \(3.09\times 10^{17}\,\mathrm{s}^{-1}\)), and explicit Markov (\(\Gamma_{\mathrm{dec}}/\omega_{\mathrm{th}} \sim 7{,}855\) at \(300\,\mathrm{K}\)) \& \(\tau_p = m\mu\) regime bounds (\nolinkurl{verify_paper062_v2.py}, 1/1 pass) & Bath spectral cutoff, non-Markovian short-time dynamics, and full microscopic bath-to-QGTCD reduction remain open \\
\textbf{Paper 063} (\nolinkurl{10.5281/zenodo.21693947}) & Four-lens chronology \& schematic analogy audit & Downgrades §§6.1--6.2 to schematic representation of bounded analogy; 4/4 null baselines match exact 4-stage CIS using exact analytical population expectations separated from finite-sample Monte Carlo checks (\nolinkurl{verify_paper063_v2.py}, 2/2 pass) & Historical priority dossier not re-adjudicated; coarse CIS alone has low discriminative specificity \\
\textbf{Paper 114} (\nolinkurl{10.5281/zenodo.23447925}) & Independent neuroscience track (SAN) & Version-aware catalog (\texttt{v1.0} Oct 17, \texttt{v1.2} Oct 20, \texttt{v1.1} supplement intentionally retained); no forced SIT scalar bridge & Calibrated comparator performance, held-out interventions, and biological data on its own track \\
\end{longtable}
\endgroup
}

\subsection{B.4 Release and Deposit Boundary}\label{b.4-release-and-deposit-boundary}

The July 2026 editions of Papers 002, 003, 007, 014, 028, 033, 034, 048, 055, 062, and 063 remain preserved immutably at their July 29, 2026 Zenodo DOIs and \texttt{pub.svgn.io} pages. All eleven November 20, 2026 corrected candidate editions and their verification scripts live in private staging pending external review and explicit public-release authorization.

\end{document}
